\documentclass[10pt]{extarticle}
\usepackage[a4paper,margin=1.5cm]{geometry}
\usepackage[T1]{fontenc}
\usepackage{amsmath,amssymb,mathtools}
\usepackage{graphicx}
\usepackage[export]{adjustbox}
\usepackage{multicol}
\usepackage{enumitem}
\usepackage{booktabs,array,tabularx}
\setlength{\columnsep}{0.7cm}
\setlength{\parindent}{0pt}
\setlength{\parskip}{3pt}
\pagestyle{empty}
\makeatletter
\newcommand{\qfigdef}[4]{\@namedef{qf@file@#1}{#2}\@namedef{qf@w@#1}{#3}\@namedef{qf@wide@#1}{#4}}
\newcommand{\qf@check}[1]{\@ifundefined{qf@file@#1}{\PackageError{score2}{Unknown figure id #1}{}}{}}
\newcommand{\qf@img}[4][]{\includegraphics[width=#2\dimexpr\@nameuse{qf@w@#4}pt\relax,max width=#3,max height=0.45\textheight,#1]{\@nameuse{qf@file@#4}}}
\newif\ifqf@full
\newif\ifqf@boxed
\newsavebox{\qf@box}
\newcommand{\qf@open}{\global\qf@fulltrue\par\global\qf@boxedtrue
  \setbox\qf@box\vbox\bgroup\hsize\textwidth\linewidth\textwidth}
\newcommand{\qf@emit}[1]{\par\egroup\global\qf@boxedfalse
  \ifdim\dimexpr\ht\qf@box+\dp\qf@box\relax<#1\textheight\noindent\usebox{\qf@box}\par\else\unvbox\qf@box\fi}
\newcommand{\qfullsplit}{\ifqf@boxed\qf@emit{0.85}\fi}
\newcommand{\qf@close}{\ifqf@boxed\qf@emit{0.45}\fi\global\qf@fullfalse}
\newcommand{\qfullbegin}{\end{multicols}\qf@open}
\newcommand{\qfullend}{\qf@close\begin{multicols}{2}}
\newcommand{\qfullstart}{\qf@open}
\newcommand{\qfulldone}{\qf@close}
\newcommand{\qfig}[2][1]{\qf@check{#2}%
  \ifqf@full\qf@blk{#1}{#2}\else
  \ifnum\@nameuse{qf@wide@#2}=1 \expandafter\@firstoftwo\else\expandafter\@secondoftwo\fi
  {\end{multicols}\begin{center}\qf@img{#1}{\textwidth}{#2}\end{center}\begin{multicols}{2}}%
  {\qf@blk{#1}{#2}}\fi}
\newcommand{\qf@blk}[2]{\par\smallskip{\centering\qf@img{#1}{\linewidth}{#2}\par}\smallskip}
\newcommand{\qfigi}[2][1]{\qf@check{#2}\qf@img[valign=c]{#1}{\linewidth}{#2}}
\makeatother
\newcommand{\papersec}[1]{\par\bigskip{\centering\Large\bfseries PAPER-#1\par}\medskip}
\newcommand{\partsec}[2]{\par\medskip\noindent{\large\bfseries PART-#1: #2}\par\smallskip}
\newcommand{\sect}[1]{\par\smallskip\noindent\textbf{#1}\par}
\newcommand{\qn}[1]{\par\smallskip\noindent\textbf{#1.}~}
\newcommand{\paraq}[2]{\par\smallskip\noindent\textbf{Paragraph for Questions #1 and #2}\par}
\newenvironment{qoptions}{\begin{itemize}[leftmargin=1.8em,itemsep=1pt,topsep=2pt,label={}]}{\end{itemize}}
\newcommand{\opt}[2]{\item[(#1)] #2}
\newcolumntype{L}[1]{>{\raggedright\arraybackslash}p{#1}}
\newcommand{\fwstart}{\end{multicols}\par\noindent\begin{minipage}{\textwidth}\small\setlength{\tabcolsep}{4pt}\renewcommand{\arraystretch}{1.5}}
\newcommand{\fwend}{\end{minipage}\par\medskip\begin{multicols}{2}}
% figure registry (generated by score2_tools.py finalize)
\qfigdef{p1_phy_s1_q01_1}{SRG (Enthusiast) - Test 13 (10 Oct 2021)-img/p1_phy_s1_q01_1.png}{145.5}{0}
\qfigdef{p1_phy_s1_q02_1}{SRG (Enthusiast) - Test 13 (10 Oct 2021)-img/p1_phy_s1_q02_1.png}{115.5}{0}
\qfigdef{p1_phy_s1_q03_1}{SRG (Enthusiast) - Test 13 (10 Oct 2021)-img/p1_phy_s1_q03_1.png}{141.4}{0}
\qfigdef{p1_phy_s1_q04_1}{SRG (Enthusiast) - Test 13 (10 Oct 2021)-img/p1_phy_s1_q04_1.png}{141.6}{0}
\qfigdef{p1_phy_s1_q05_1}{SRG (Enthusiast) - Test 13 (10 Oct 2021)-img/p1_phy_s1_q05_1.png}{116.9}{0}
\qfigdef{p1_phy_s2_q06_1}{SRG (Enthusiast) - Test 13 (10 Oct 2021)-img/p1_phy_s2_q06_1.png}{124.4}{0}
\qfigdef{p1_phy_s2_q07_1}{SRG (Enthusiast) - Test 13 (10 Oct 2021)-img/p1_phy_s2_q07_1.png}{140.5}{0}
\qfigdef{p1_phy_s2_q08_1}{SRG (Enthusiast) - Test 13 (10 Oct 2021)-img/p1_phy_s2_q08_1.png}{124.2}{0}
\qfigdef{p1_phy_s2_q09_A}{SRG (Enthusiast) - Test 13 (10 Oct 2021)-img/p1_phy_s2_q09_A.png}{81.2}{0}
\qfigdef{p1_phy_s2_q09_2}{SRG (Enthusiast) - Test 13 (10 Oct 2021)-img/p1_phy_s2_q09_2.png}{94.0}{0}
\qfigdef{p1_phy_s2_q09_C}{SRG (Enthusiast) - Test 13 (10 Oct 2021)-img/p1_phy_s2_q09_C.png}{75.1}{0}
\qfigdef{p1_phy_s2_q09_D}{SRG (Enthusiast) - Test 13 (10 Oct 2021)-img/p1_phy_s2_q09_D.png}{82.0}{0}
\qfigdef{p1_phy_s2_q10_1}{SRG (Enthusiast) - Test 13 (10 Oct 2021)-img/p1_phy_s2_q10_1.png}{171.5}{0}
\qfigdef{p1_phy_s2_q11_1}{SRG (Enthusiast) - Test 13 (10 Oct 2021)-img/p1_phy_s2_q11_1.png}{134.0}{0}
\qfigdef{p1_phy_s2_q12_1}{SRG (Enthusiast) - Test 13 (10 Oct 2021)-img/p1_phy_s2_q12_1.png}{87.7}{0}
\qfigdef{p1_phy_s3_q14_1}{SRG (Enthusiast) - Test 13 (10 Oct 2021)-img/p1_phy_s3_q14_1.png}{97.1}{0}
\qfigdef{p1_phy_s3_q15_1}{SRG (Enthusiast) - Test 13 (10 Oct 2021)-img/p1_phy_s3_q15_1.png}{211.6}{0}
\qfigdef{p1_phy_s3_q16_1}{SRG (Enthusiast) - Test 13 (10 Oct 2021)-img/p1_phy_s3_q16_1.png}{206.6}{0}
\qfigdef{p1_phy_s3_q17_1}{SRG (Enthusiast) - Test 13 (10 Oct 2021)-img/p1_phy_s3_q17_1.png}{108.5}{0}
\qfigdef{p1_phy_s3_q18_1}{SRG (Enthusiast) - Test 13 (10 Oct 2021)-img/p1_phy_s3_q18_1.png}{225.8}{0}
\qfigdef{p1_chem_s1_q01_1}{SRG (Enthusiast) - Test 13 (10 Oct 2021)-img/p1_chem_s1_q01_1.png}{62.0}{0}
\qfigdef{p1_chem_s1_q01_A}{SRG (Enthusiast) - Test 13 (10 Oct 2021)-img/p1_chem_s1_q01_A.png}{50.4}{0}
\qfigdef{p1_chem_s1_q01_B}{SRG (Enthusiast) - Test 13 (10 Oct 2021)-img/p1_chem_s1_q01_B.png}{50.8}{0}
\qfigdef{p1_chem_s1_q01_C}{SRG (Enthusiast) - Test 13 (10 Oct 2021)-img/p1_chem_s1_q01_C.png}{68.5}{0}
\qfigdef{p1_chem_s1_q01_D}{SRG (Enthusiast) - Test 13 (10 Oct 2021)-img/p1_chem_s1_q01_D.png}{65.9}{0}
\qfigdef{p1_chem_s1_q02_1}{SRG (Enthusiast) - Test 13 (10 Oct 2021)-img/p1_chem_s1_q02_1.png}{205.1}{0}
\qfigdef{p1_chem_s1_q02_A}{SRG (Enthusiast) - Test 13 (10 Oct 2021)-img/p1_chem_s1_q02_A.png}{63.9}{0}
\qfigdef{p1_chem_s1_q02_B}{SRG (Enthusiast) - Test 13 (10 Oct 2021)-img/p1_chem_s1_q02_B.png}{81.4}{0}
\qfigdef{p1_chem_s1_q02_C}{SRG (Enthusiast) - Test 13 (10 Oct 2021)-img/p1_chem_s1_q02_C.png}{63.9}{0}
\qfigdef{p1_chem_s1_q02_D}{SRG (Enthusiast) - Test 13 (10 Oct 2021)-img/p1_chem_s1_q02_D.png}{82.3}{0}
\qfigdef{p1_chem_s1_q03_1}{SRG (Enthusiast) - Test 13 (10 Oct 2021)-img/p1_chem_s1_q03_1.png}{164.5}{0}
\qfigdef{p1_chem_s1_q03_A}{SRG (Enthusiast) - Test 13 (10 Oct 2021)-img/p1_chem_s1_q03_A.png}{69.4}{0}
\qfigdef{p1_chem_s1_q03_B}{SRG (Enthusiast) - Test 13 (10 Oct 2021)-img/p1_chem_s1_q03_B.png}{55.0}{0}
\qfigdef{p1_chem_s1_q03_C}{SRG (Enthusiast) - Test 13 (10 Oct 2021)-img/p1_chem_s1_q03_C.png}{45.2}{0}
\qfigdef{p1_chem_s1_q03_D}{SRG (Enthusiast) - Test 13 (10 Oct 2021)-img/p1_chem_s1_q03_D.png}{46.5}{0}
\qfigdef{p1_chem_s1_q04_1}{SRG (Enthusiast) - Test 13 (10 Oct 2021)-img/p1_chem_s1_q04_1.png}{106.5}{0}
\qfigdef{p1_chem_s1_q04_A}{SRG (Enthusiast) - Test 13 (10 Oct 2021)-img/p1_chem_s1_q04_A.png}{61.5}{0}
\qfigdef{p1_chem_s1_q04_B}{SRG (Enthusiast) - Test 13 (10 Oct 2021)-img/p1_chem_s1_q04_B.png}{85.3}{0}
\qfigdef{p1_chem_s1_q04_C}{SRG (Enthusiast) - Test 13 (10 Oct 2021)-img/p1_chem_s1_q04_C.png}{134.0}{0}
\qfigdef{p1_chem_s1_q04_D}{SRG (Enthusiast) - Test 13 (10 Oct 2021)-img/p1_chem_s1_q04_D.png}{60.9}{0}
\qfigdef{p1_chem_s2_q06_A}{SRG (Enthusiast) - Test 13 (10 Oct 2021)-img/p1_chem_s2_q06_A.png}{174.7}{0}
\qfigdef{p1_chem_s2_q06_B}{SRG (Enthusiast) - Test 13 (10 Oct 2021)-img/p1_chem_s2_q06_B.png}{124.4}{0}
\qfigdef{p1_chem_s2_q06_C}{SRG (Enthusiast) - Test 13 (10 Oct 2021)-img/p1_chem_s2_q06_C.png}{186.1}{0}
\qfigdef{p1_chem_s2_q06_D}{SRG (Enthusiast) - Test 13 (10 Oct 2021)-img/p1_chem_s2_q06_D.png}{93.2}{0}
\qfigdef{p1_chem_s2_q07_A}{SRG (Enthusiast) - Test 13 (10 Oct 2021)-img/p1_chem_s2_q07_A.png}{241.5}{0}
\qfigdef{p1_chem_s2_q07_3}{SRG (Enthusiast) - Test 13 (10 Oct 2021)-img/p1_chem_s2_q07_3.png}{235.8}{0}
\qfigdef{p1_chem_s2_q07_4}{SRG (Enthusiast) - Test 13 (10 Oct 2021)-img/p1_chem_s2_q07_4.png}{166.5}{0}
\qfigdef{p1_chem_s2_q07_5}{SRG (Enthusiast) - Test 13 (10 Oct 2021)-img/p1_chem_s2_q07_5.png}{208.2}{0}
\qfigdef{p1_chem_s2_q08_A}{SRG (Enthusiast) - Test 13 (10 Oct 2021)-img/p1_chem_s2_q08_A.png}{40.5}{0}
\qfigdef{p1_chem_s2_q08_B}{SRG (Enthusiast) - Test 13 (10 Oct 2021)-img/p1_chem_s2_q08_B.png}{73.5}{0}
\qfigdef{p1_chem_s2_q08_C}{SRG (Enthusiast) - Test 13 (10 Oct 2021)-img/p1_chem_s2_q08_C.png}{34.7}{0}
\qfigdef{p1_chem_s2_q08_D}{SRG (Enthusiast) - Test 13 (10 Oct 2021)-img/p1_chem_s2_q08_D.png}{68.3}{0}
\qfigdef{p1_chem_s2_q09_1}{SRG (Enthusiast) - Test 13 (10 Oct 2021)-img/p1_chem_s2_q09_1.png}{115.2}{0}
\qfigdef{p1_chem_s2_q13_1}{SRG (Enthusiast) - Test 13 (10 Oct 2021)-img/p1_chem_s2_q13_1.png}{188.7}{0}
\qfigdef{p1_chem_s2_q13_A}{SRG (Enthusiast) - Test 13 (10 Oct 2021)-img/p1_chem_s2_q13_A.png}{63.7}{0}
\qfigdef{p1_chem_s2_q13_B}{SRG (Enthusiast) - Test 13 (10 Oct 2021)-img/p1_chem_s2_q13_B.png}{59.5}{0}
\qfigdef{p1_chem_s2_q13_C}{SRG (Enthusiast) - Test 13 (10 Oct 2021)-img/p1_chem_s2_q13_C.png}{63.7}{0}
\qfigdef{p1_chem_s2_q13_D}{SRG (Enthusiast) - Test 13 (10 Oct 2021)-img/p1_chem_s2_q13_D.png}{61.3}{0}
\qfigdef{p1_chem_s3_q15_1}{SRG (Enthusiast) - Test 13 (10 Oct 2021)-img/p1_chem_s3_q15_1.png}{364.2}{1}
\qfigdef{p1_chem_s3_q17_9}{SRG (Enthusiast) - Test 13 (10 Oct 2021)-img/p1_chem_s3_q17_9.png}{412.2}{1}
\qfigdef{p1_chem_s3_q18_1}{SRG (Enthusiast) - Test 13 (10 Oct 2021)-img/p1_chem_s3_q18_1.png}{389.3}{1}
\qfigdef{p2_phy_s1_q01_1}{SRG (Enthusiast) - Test 13 (10 Oct 2021)-img/p2_phy_s1_q01_1.png}{178.0}{0}
\qfigdef{p2_phy_s1_q02_1}{SRG (Enthusiast) - Test 13 (10 Oct 2021)-img/p2_phy_s1_q02_1.png}{68.7}{0}
\qfigdef{p2_phy_s1_q03_1}{SRG (Enthusiast) - Test 13 (10 Oct 2021)-img/p2_phy_s1_q03_1.png}{116.7}{0}
\qfigdef{p2_phy_s1_q04_1}{SRG (Enthusiast) - Test 13 (10 Oct 2021)-img/p2_phy_s1_q04_1.png}{130.7}{0}
\qfigdef{p2_phy_s1_q05_1}{SRG (Enthusiast) - Test 13 (10 Oct 2021)-img/p2_phy_s1_q05_1.png}{175.2}{0}
\qfigdef{p2_phy_s1_q05_2}{SRG (Enthusiast) - Test 13 (10 Oct 2021)-img/p2_phy_s1_q05_2.png}{122.8}{0}
\qfigdef{p2_phy_s1_q07_1}{SRG (Enthusiast) - Test 13 (10 Oct 2021)-img/p2_phy_s1_q07_1.png}{116.7}{0}
\qfigdef{p2_phy_s1_q08_1}{SRG (Enthusiast) - Test 13 (10 Oct 2021)-img/p2_phy_s1_q08_1.png}{146.6}{0}
\qfigdef{p2_phy_s2_q09_A}{SRG (Enthusiast) - Test 13 (10 Oct 2021)-img/p2_phy_s2_q09_A.png}{63.5}{0}
\qfigdef{p2_phy_s2_q09_B}{SRG (Enthusiast) - Test 13 (10 Oct 2021)-img/p2_phy_s2_q09_B.png}{87.9}{0}
\qfigdef{p2_phy_s2_q09_3}{SRG (Enthusiast) - Test 13 (10 Oct 2021)-img/p2_phy_s2_q09_3.png}{92.5}{0}
\qfigdef{p2_phy_s2_q09_D}{SRG (Enthusiast) - Test 13 (10 Oct 2021)-img/p2_phy_s2_q09_D.png}{117.4}{0}
\qfigdef{p2_phy_s2_q10_A}{SRG (Enthusiast) - Test 13 (10 Oct 2021)-img/p2_phy_s2_q10_A.png}{101.0}{0}
\qfigdef{p2_phy_s2_q10_2}{SRG (Enthusiast) - Test 13 (10 Oct 2021)-img/p2_phy_s2_q10_2.png}{140.5}{0}
\qfigdef{p2_phy_s2_q10_B}{SRG (Enthusiast) - Test 13 (10 Oct 2021)-img/p2_phy_s2_q10_B.png}{101.0}{0}
\qfigdef{p2_phy_s2_q10_3}{SRG (Enthusiast) - Test 13 (10 Oct 2021)-img/p2_phy_s2_q10_3.png}{140.5}{0}
\qfigdef{p2_phy_s2_q10_C}{SRG (Enthusiast) - Test 13 (10 Oct 2021)-img/p2_phy_s2_q10_C.png}{101.0}{0}
\qfigdef{p2_phy_s2_q10_5}{SRG (Enthusiast) - Test 13 (10 Oct 2021)-img/p2_phy_s2_q10_5.png}{140.5}{0}
\qfigdef{p2_phy_s2_q10_D}{SRG (Enthusiast) - Test 13 (10 Oct 2021)-img/p2_phy_s2_q10_D.png}{101.0}{0}
\qfigdef{p2_phy_s2_q10_7}{SRG (Enthusiast) - Test 13 (10 Oct 2021)-img/p2_phy_s2_q10_7.png}{140.5}{0}
\qfigdef{p2_phy_s2_q10_9}{SRG (Enthusiast) - Test 13 (10 Oct 2021)-img/p2_phy_s2_q10_9.png}{139.8}{0}
\qfigdef{p2_phy_s3_q11_1}{SRG (Enthusiast) - Test 13 (10 Oct 2021)-img/p2_phy_s3_q11_1.png}{98.4}{0}
\qfigdef{p2_phy_s3_q12_1}{SRG (Enthusiast) - Test 13 (10 Oct 2021)-img/p2_phy_s3_q12_1.png}{105.8}{0}
\qfigdef{p2_phy_s3_q15_1}{SRG (Enthusiast) - Test 13 (10 Oct 2021)-img/p2_phy_s3_q15_1.png}{66.1}{0}
\qfigdef{p2_phy_s3_q16_2}{SRG (Enthusiast) - Test 13 (10 Oct 2021)-img/p2_phy_s3_q16_2.png}{195.7}{0}
\qfigdef{p2_phy_s3_q17_1}{SRG (Enthusiast) - Test 13 (10 Oct 2021)-img/p2_phy_s3_q17_1.png}{65.6}{0}
\qfigdef{p2_phy_s3_q18_1}{SRG (Enthusiast) - Test 13 (10 Oct 2021)-img/p2_phy_s3_q18_1.png}{134.2}{0}
\qfigdef{p2_chem_s1_q01_1}{SRG (Enthusiast) - Test 13 (10 Oct 2021)-img/p2_chem_s1_q01_1.png}{120.9}{0}
\qfigdef{p2_chem_s1_q02_1}{SRG (Enthusiast) - Test 13 (10 Oct 2021)-img/p2_chem_s1_q02_1.png}{42.1}{0}
\qfigdef{p2_chem_s1_q02_A}{SRG (Enthusiast) - Test 13 (10 Oct 2021)-img/p2_chem_s1_q02_A.png}{104.3}{0}
\qfigdef{p2_chem_s1_q02_B}{SRG (Enthusiast) - Test 13 (10 Oct 2021)-img/p2_chem_s1_q02_B.png}{100.4}{0}
\qfigdef{p2_chem_s1_q02_C}{SRG (Enthusiast) - Test 13 (10 Oct 2021)-img/p2_chem_s1_q02_C.png}{101.0}{0}
\qfigdef{p2_chem_s1_q02_D}{SRG (Enthusiast) - Test 13 (10 Oct 2021)-img/p2_chem_s1_q02_D.png}{100.4}{0}
\qfigdef{p2_chem_s1_q03_1}{SRG (Enthusiast) - Test 13 (10 Oct 2021)-img/p2_chem_s1_q03_1.png}{172.2}{0}
\qfigdef{p2_chem_s1_q03_A}{SRG (Enthusiast) - Test 13 (10 Oct 2021)-img/p2_chem_s1_q03_A.png}{79.6}{0}
\qfigdef{p2_chem_s1_q03_B}{SRG (Enthusiast) - Test 13 (10 Oct 2021)-img/p2_chem_s1_q03_B.png}{141.2}{0}
\qfigdef{p2_chem_s1_q03_D}{SRG (Enthusiast) - Test 13 (10 Oct 2021)-img/p2_chem_s1_q03_D.png}{146.8}{0}
\qfigdef{p2_chem_s1_q04_A}{SRG (Enthusiast) - Test 13 (10 Oct 2021)-img/p2_chem_s1_q04_A.png}{258.4}{1}
\qfigdef{p2_chem_s1_q04_C}{SRG (Enthusiast) - Test 13 (10 Oct 2021)-img/p2_chem_s1_q04_C.png}{133.1}{0}
\qfigdef{p2_chem_s1_q04_D}{SRG (Enthusiast) - Test 13 (10 Oct 2021)-img/p2_chem_s1_q04_D.png}{237.2}{0}
\qfigdef{p2_chem_s1_q05_1}{SRG (Enthusiast) - Test 13 (10 Oct 2021)-img/p2_chem_s1_q05_1.png}{55.8}{0}
\qfigdef{p2_chem_s1_q05_A}{SRG (Enthusiast) - Test 13 (10 Oct 2021)-img/p2_chem_s1_q05_A.png}{84.6}{0}
\qfigdef{p2_chem_s1_q05_B}{SRG (Enthusiast) - Test 13 (10 Oct 2021)-img/p2_chem_s1_q05_B.png}{76.2}{0}
\qfigdef{p2_chem_s1_q05_C}{SRG (Enthusiast) - Test 13 (10 Oct 2021)-img/p2_chem_s1_q05_C.png}{67.6}{0}
\qfigdef{p2_chem_s1_q05_D}{SRG (Enthusiast) - Test 13 (10 Oct 2021)-img/p2_chem_s1_q05_D.png}{75.3}{0}
\qfigdef{p2_chem_s1_q07_1}{SRG (Enthusiast) - Test 13 (10 Oct 2021)-img/p2_chem_s1_q07_1.png}{40.4}{0}
\qfigdef{p2_chem_s2_q09_A}{SRG (Enthusiast) - Test 13 (10 Oct 2021)-img/p2_chem_s2_q09_A.png}{81.4}{0}
\qfigdef{p2_chem_s2_q09_B}{SRG (Enthusiast) - Test 13 (10 Oct 2021)-img/p2_chem_s2_q09_B.png}{74.8}{0}
\qfigdef{p2_chem_s2_q09_C}{SRG (Enthusiast) - Test 13 (10 Oct 2021)-img/p2_chem_s2_q09_C.png}{68.9}{0}
\qfigdef{p2_chem_s2_q09_D}{SRG (Enthusiast) - Test 13 (10 Oct 2021)-img/p2_chem_s2_q09_D.png}{135.7}{0}
\qfigdef{p2_chem_s3_q11_1}{SRG (Enthusiast) - Test 13 (10 Oct 2021)-img/p2_chem_s3_q11_1.png}{232.8}{0}
\qfigdef{p2_chem_s3_q13_10}{SRG (Enthusiast) - Test 13 (10 Oct 2021)-img/p2_chem_s3_q13_10.png}{388.4}{1}
\qfigdef{p2_chem_s3_q15_1}{SRG (Enthusiast) - Test 13 (10 Oct 2021)-img/p2_chem_s3_q15_1.png}{113.5}{0}
\qfigdef{p2_chem_s3_q15_2}{SRG (Enthusiast) - Test 13 (10 Oct 2021)-img/p2_chem_s3_q15_2.png}{128.1}{0}
\qfigdef{p2_chem_s3_q15_3}{SRG (Enthusiast) - Test 13 (10 Oct 2021)-img/p2_chem_s3_q15_3.png}{126.4}{0}
\qfigdef{p2_chem_s3_q16_9}{SRG (Enthusiast) - Test 13 (10 Oct 2021)-img/p2_chem_s3_q16_9.png}{395.5}{1}
\qfigdef{p2_chem_s3_q18_1}{SRG (Enthusiast) - Test 13 (10 Oct 2021)-img/p2_chem_s3_q18_1.png}{99.5}{0}
\raggedcolumns\makeatletter\@beginparpenalty=10000\makeatother

\begin{document}
\begin{center}{\Large\bfseries SRG (Enthusiast) --- Test 13 (10 Oct 2021)}\end{center}
\begin{multicols}{2}
\papersec{1}
\partsec{1}{PHYSICS}
\sect{SECTION-I (i)}
\qn{1} In the adjacent figure, a block of mass $m$ is rigidly attached to two identical spring of stiffness $K$ each. The other ends of the spring are connected to the fixed wall. When the block is in equilibrium, length of each spring is $b$, which is greater than the natural length $l$ of the spring. The Time -- period of the oscillation of the block if it is displaced by small distance $X$ perpendicular to the length of the spring and released. (neglect gravity)
\qfig{p1_phy_s1_q01_1}
\begin{qoptions}
\opt{A}{$2\pi\sqrt{\dfrac{mb}{k(b-l)}}$}
\opt{B}{$2\pi\sqrt{\dfrac{mb}{2k(b-l)}}$}
\opt{C}{$2\pi\sqrt{\dfrac{m(b-l)}{kb}}$}
\opt{D}{$2\pi\sqrt{\dfrac{m(b-l)}{2kb}}$}
\end{qoptions}
\qn{2} The semi-circular cylindrical shell of radius $R$ with small but uniform wall thickness is set into small oscillation on the horizontal surface. If no slipping occurs, then the expression for the period $T$ of each complete oscillation is :
\qfig{p1_phy_s1_q02_1}
\begin{qoptions}
\opt{A}{$T=2\pi\sqrt{\dfrac{(\pi-2)R}{g}}$}
\opt{B}{$T=2\pi\sqrt{\dfrac{(\pi-4)R}{g}}$}
\opt{C}{$T=2\pi\sqrt{\dfrac{2R}{g}}$}
\opt{D}{$T=2\pi\sqrt{\dfrac{(\pi-6)R}{g}}$}
\end{qoptions}
\qn{3} The instantaneous potential difference between points $A$ and $B$ is :
\qfig{p1_phy_s1_q03_1}
\begin{qoptions}
\opt{A}{$8\sin\!\left(50\pi t+37\dfrac{\pi}{180}\right)$}
\opt{B}{$8\sin\!\left(50\pi t-37\dfrac{\pi}{180}\right)$}
\opt{C}{$10\sin(50\pi t)$}
\opt{D}{$10\cos(50\pi t)$}
\end{qoptions}
\qn{4} A square loop of side `$a$' is placed in $x$-$y$ plane as shown in figure. In this region there is non-uniform time dependent magnetic field $\vec B=\left(cy^3t^2\right)\hat k$. [where $t$ is time and $c$ is constant] then magnitude of emf induced in loop is
\qfig{p1_phy_s1_q04_1}
\begin{qoptions}
\opt{A}{Zero}
\opt{B}{$\dfrac{ca^4t^2}{5}$}
\opt{C}{$\dfrac{ca^5t}{2}$}
\opt{D}{Data is insufficient}
\end{qoptions}
\qn{5} Three identical large plates are fixed at separation of $d$ from each other as shown. The area of each plate is $A$. Plate 1 is given charge $+Q$ while plates 2 and 3 are neutral and are connected to each other through coil of inductances $L$ and switch $S$. If resistance of all connected wires is neglected the maximum current flow through coil after closing switch is $\left(C=\varepsilon_0A/d\right)$ (neglect fringe effect)
\qfig{p1_phy_s1_q05_1}
\begin{qoptions}
\opt{A}{$\dfrac{Q_0}{\sqrt{LC}}$}
\opt{B}{$\dfrac{Q_0}{\sqrt{2LC}}$}
\opt{C}{$\dfrac{2Q_0}{\sqrt{LC}}$}
\opt{D}{$\dfrac{Q_0}{2\sqrt{LC}}$}
\end{qoptions}
\sect{SECTION-I (ii)}
\qn{6} A simple pendulum consists of a bob of mass $m$ and a light string of length $l$ as shown in the figure. Another identical ball moving with the small velocity $v_0$ collides with the pendulum's bob and sticks to it. For this new pendulum of mass $2m$, mark out the correct statement(s)
\qfig{p1_phy_s2_q06_1}
\begin{qoptions}
\opt{A}{Time period of the pendulum is $2\pi\sqrt{\dfrac lg}$}
\opt{B}{The equation of motion for this pendulum is $\theta=\dfrac{v_0}{2\sqrt{gl}}\sin\!\left[\sqrt{\dfrac gl}\,t\right]$}
\opt{C}{The equation of motion for this pendulum is $\theta=\dfrac{v_0}{2\sqrt{gl}}\cos\!\left[\sqrt{\dfrac gl}\,t\right]$}
\opt{D}{Time period of the pendulum is $2\pi\sqrt{\dfrac{2l}{g}}$}
\end{qoptions}
\qn{7} A wire shaped as a semicircle of radius $a$, is rotating about an axis $PQ$ with a constant angular velocity $\omega=\dfrac{1}{\sqrt{LC}}$, with the help of an external agent. A uniform magnetic field $B$ exists in space and is directed into the plane of the figure. (circuit part remains at rest) (left part is at rest)
\qfig{p1_phy_s2_q07_1}
\begin{qoptions}
\opt{A}{The rms value of current in the circuit is $\dfrac{\pi Ba^2}{R\sqrt{2LC}}$}
\opt{B}{The rms value of current in the circuit is $\dfrac{\pi Ba^2}{2R\sqrt{2LC}}$}
\opt{C}{The maximum energy stored in the capacitor is $\dfrac{\pi^2B^2a^4}{8R^2C}$}
\opt{D}{The maximum power delivered by the external agent is $\dfrac{\pi^2B^2a^4}{4LCR}$}
\end{qoptions}
\qn{8} In the figure shown the key is switched on at $t=0$. Let $I_1$ and $I_2$ be the currents through inductors having self inductances $L_1$ \& $L_2$ at any time $t$ respectively. The magnetic energy stored in the inductors 1 and 2 be $U_1$ and $U_2$. Then $\dfrac{U_1}{U_2}$ at any instant of time is :
\qfig{p1_phy_s2_q08_1}
\begin{qoptions}
\opt{A}{$\dfrac{L_1}{L_2}$}
\opt{B}{$\dfrac{L_2}{L_1}$}
\opt{C}{$\dfrac{I_1}{I_2}$}
\opt{D}{$\dfrac{I_2}{I_1}$}
\end{qoptions}
\qn{9} A long straight slider is moved with constant velocity $\vec v=v_0\hat i$ parallel to earth surface in the $x$-$y$ plane in a uniform magnetic field $\vec B=B_0\hat k$ on a resistanceless rails in the form of curve $y^2=4x$. Here emf induced in the loop, current, power, resistance of slider per unit length and time are represented by $e,i,p,\lambda_0$ and $t$ respectively. If at $t=0$ slider starts from $x=0$. Then which of the following option(s) is / are correct.
\begin{qoptions}
\opt{A}{\qfigi{p1_phy_s2_q09_A}}
\opt{B}{\qfigi{p1_phy_s2_q09_2}}
\opt{C}{\qfigi{p1_phy_s2_q09_C}}
\opt{D}{\qfigi{p1_phy_s2_q09_D}}
\end{qoptions}
\qn{10} In the figure shows $R=100\,\Omega$, $L=\dfrac2\pi$ H and $C=\dfrac8\pi\,\mu$F are connected in series with an a.c source of 200 volt and frequency $f$. $V_1$ and $V_2$ are two hot-wire voltmeters. If the reading of $V_1$ and $V_2$ are same, then
\qfig{p1_phy_s2_q10_1}
\begin{qoptions}
\opt{A}{$f=125$ Hz}
\opt{B}{$f=250\pi$ Hz}
\opt{C}{Current through $R$ is 2A}
\opt{D}{$V_1=V_2=1000$ volt}
\end{qoptions}
\qn{11} A simple pendulum initially oscillating simple harmonically with angular amplitude $\alpha$ and time period $T_0$ is symmetrically confined between two rigid fixed planes $P$ and $Q$ making angle `$\alpha$' with each other as shown in figure. [Assume `$\alpha$' is small]
\qfig{p1_phy_s2_q11_1}
\begin{qoptions}
\opt{A}{If collision at both walls are elastic then time period is $\dfrac{T_0}{2}$}
\opt{B}{If collision at both walls are inelastic then time period is $T_0$}
\opt{C}{If collision at one wall is elastic and at other is inelastic, then 1 time period is $T_0$}
\opt{D}{If collision at one wall is elastic and at other is inelastic, then 1 time period will be less than $T_0$.}
\end{qoptions}
\qn{12} In the cylindrical region of radius $R=2$ m, there exists a time varying magnetic field $B$ such that $\dfrac{dB}{dt}=2$ Tesla/sec. A charged particle having charge $q=2$ C is placed at the point $P$ at a distance $d$ from its center $O$. Now, the particle is moved in the direction perpendicular to line $OP$ by an external agent up to infinity so that there is no gain in kinetic energy of charged particle. Then
\qfig{p1_phy_s2_q12_1}
\begin{qoptions}
\opt{A}{Work done by external agent is $4\pi$ Joule when $d=4$ m}
\opt{B}{Work done by external agent is $4\pi$ Joule when $d=8$ m}
\opt{C}{Work done by external agent is independent of $d$}
\opt{D}{Work done by external agent is positive}
\end{qoptions}
\qn{13} A uniform conducting ring of mass $m=2$ kg, radius $r=2$ m and resistance $8\,\Omega$ is kept on smooth, horizontal surface. A time varying magnetic field $\vec B=\left(\hat i+t^2\hat j\right)$ Tesla is present in the region, where $t$ is time in second and take vertical as $y$-axis. (Take $\pi^2=10$). Then
\begin{qoptions}
\opt{A}{Time when ring starts toppling is 1 sec}
\opt{B}{Time when ring starts toppling is 3/4 sec}
\opt{C}{Heat generated through the ring till the instant when the ring start toppling is 40/3 Joule}
\opt{D}{Heat generated through the ring till the instant when the ring start toppling is 80/3 Joule}
\end{qoptions}
\sect{SECTION-II}
\qn{14} A meter stick connected with disc of same mass which is joined at one end of stick and disc can freely rotates about its joint. System swings in a vertical plane about a fixed horizontal axis passing through its one end, undergoes small oscillation of frequency `$f_0$'. If disc is removed from the stick, then its new frequency of small oscillation is $\sqrt{\dfrac{35}{k}}f_0$, find the value of $k$. (take $R=l/2$)
\qfig{p1_phy_s3_q14_1}
\qn{15} In the circuit the key $(K)$ is closed at $t=0$. Find the current in ampere through the key $(K)$ at the instant $t=10^{-3}\ln2$ sec
\qfig{p1_phy_s3_q15_1}
\qn{16} Consider a network shown in figure. The current supplied by the AC voltage source is $n$ amperes. What is the value of $n$.
\qfig{p1_phy_s3_q16_1}
\qn{17} A conducting rod of mass $m=0.3$ kg and length $l=4$ m can slide without friction on two parallel conducting rails. The conducting rails are connected via an inductance $L=3$ mH. This system is placed in a region containing uniform magnetic field $B=1$ T pointing into the plane. If the rod is given an initial velocity $v_0=2$ m/s, its oscillates with an amplitude $A$ cm. Find the value of $2A$.
\qfig{p1_phy_s3_q17_1}
\qn{18} A spring block system is placed on a horizontal rough surface as shown in the figure. The block is given velocity when the spring is in natural length. The total distance travelled by the block before it finally comes to rest is
\qfig{p1_phy_s3_q18_1}
\partsec{2}{CHEMISTRY}
\sect{SECTION-I (i)}
\qn{1} \qfigi{p1_chem_s1_q01_1} $\xrightarrow{\Delta}$ major product is :
\begin{qoptions}
\opt{A}{\qfigi{p1_chem_s1_q01_A}}
\opt{B}{\qfigi{p1_chem_s1_q01_B}}
\opt{C}{\qfigi{p1_chem_s1_q01_C}}
\opt{D}{\qfigi{p1_chem_s1_q01_D}}
\end{qoptions}
\qn{2} \qfigi{p1_chem_s1_q02_1}

`$B$' is :
\begin{qoptions}
\opt{A}{\qfigi{p1_chem_s1_q02_A}}
\opt{B}{\qfigi{p1_chem_s1_q02_B}}
\opt{C}{\qfigi{p1_chem_s1_q02_C}}
\opt{D}{\qfigi{p1_chem_s1_q02_D}}
\end{qoptions}
\qn{3} \qfigi{p1_chem_s1_q03_1}

`$X$' obtained :
\begin{qoptions}
\opt{A}{\qfigi{p1_chem_s1_q03_A}}
\opt{B}{\qfigi{p1_chem_s1_q03_B}}
\opt{C}{\qfigi{p1_chem_s1_q03_C}}
\opt{D}{\qfigi{p1_chem_s1_q03_D}}
\end{qoptions}
\qn{4} \qfigi{p1_chem_s1_q04_1}
\begin{qoptions}
\opt{A}{\qfigi{p1_chem_s1_q04_A}}
\opt{B}{\qfigi{p1_chem_s1_q04_B}}
\opt{C}{\qfigi{p1_chem_s1_q04_C}}
\opt{D}{\qfigi{p1_chem_s1_q04_D}}
\end{qoptions}
\qn{5} When 1 mol of a mono atomic ideal gas at $T$ K undergoes adiabatic change under a constant external pressure of 1 atm, changes volume from 1 L to 2 L. The final temperature (in K) would be
\begin{qoptions}
\opt{A}{$\dfrac{T}{2^{2/3}}$}
\opt{B}{$T+\dfrac{2}{3\times0.0821}$}
\opt{C}{$T$}
\opt{D}{$T-\dfrac{2}{3\times0.0821}$}
\end{qoptions}
\sect{SECTION-I (ii)}
\qn{6} Which of the following is / are correctly matched
\begin{qoptions}
\opt{A}{\qfigi{p1_chem_s2_q06_A}

No.\,of stereomers of $(P)=3$}
\opt{B}{\qfigi{p1_chem_s2_q06_B}

Then, 1 mole of $(Q)$ liberates 44.8 lit of $\mathrm{CH_4}$ gas with excess of $\mathrm{CH_3MgBr}$ / ether , at N.T.P}
\opt{C}{\qfigi{p1_chem_s2_q06_C}}
\opt{D}{\qfigi{p1_chem_s2_q06_D} gives yellow precipitation with $\mathrm{NaOBr_{(aq)}}$}
\end{qoptions}
\qn{7} Select option(s) in which correct major product is not given :
\begin{qoptions}
\opt{A}{\qfigi{p1_chem_s2_q07_A}}
\opt{B}{\qfigi{p1_chem_s2_q07_3}}
\opt{C}{\qfigi{p1_chem_s2_q07_4}}
\opt{D}{\qfigi{p1_chem_s2_q07_5}}
\end{qoptions}
\qn{8} Amongst the following compound(s) which is / are soluble in aqueous NaOH ?
\begin{qoptions}
\opt{A}{\qfigi{p1_chem_s2_q08_A}}
\opt{B}{\qfigi{p1_chem_s2_q08_B}}
\opt{C}{\qfigi{p1_chem_s2_q08_C}}
\opt{D}{\qfigi{p1_chem_s2_q08_D}}
\end{qoptions}
\qn{9} For the given conversion :

\qfigi{p1_chem_s2_q09_1}

Suitable options are :
\begin{qoptions}
\opt{A}{Reduction by $\mathrm{LiAlH_4}$}
\opt{B}{Reaction with $\mathrm{SOCl_2}$ followed by $\mathrm{H_2/Pd-BaSO_4}$}
\opt{C}{Catalytic reduction by $\mathrm{H_2/Ni}$}
\opt{D}{Reaction with $\mathrm{PCl_5}$ followed by $\mathrm{DIBAL-H/-78^\circ C}$}
\end{qoptions}
\qn{10} Compound(s) without carboxylic acid group in it are :
\begin{qoptions}
\opt{A}{Barbituric acid}
\opt{B}{Carbolic acid}
\opt{C}{Picric acid}
\opt{D}{Cyanic acid}
\end{qoptions}
\qn{11} The normal boiling point of a liquid X is 400 K. $\Delta H_{\mathrm{vap}}$ at normal boiling point is 40 kJ/mol. Select correct statement(s) :
\begin{qoptions}
\opt{A}{$\Delta S_{\mathrm{vaporisation}}<100$ J/mol.K at 400 K and 2 atm}
\opt{B}{$\Delta S_{\mathrm{vaporisation}}<10$ J/mol.K at 400 K and 1 atm}
\opt{C}{$\Delta G_{\mathrm{vaporisation}}<0$ at 410 K and 1 atm}
\opt{D}{$\Delta U=43.32$ kJ/mol.K at 400 K and 1 atm}
\end{qoptions}
\qn{12} For an isolated system, the entropy :
\begin{qoptions}
\opt{A}{Either increases or remains constant}
\opt{B}{Either decreases or remains constant}
\opt{C}{Can never decrease}
\opt{D}{Can never increase}
\end{qoptions}
\qn{13} \qfigi{p1_chem_s2_q13_1}

Product (s) is / are :
\begin{qoptions}
\opt{A}{\qfigi{p1_chem_s2_q13_A}}
\opt{B}{\qfigi{p1_chem_s2_q13_B}}
\opt{C}{\qfigi{p1_chem_s2_q13_C}}
\opt{D}{\qfigi{p1_chem_s2_q13_D}}
\end{qoptions}
\sect{SECTION-II}
\qn{14} The number of compounds (molar mass $=112$) with six membered ring, show `+' ve Brady's test and give optically active product with $\mathrm{NaBH_4}=$ `X' (Including stereomers) then the value of `12X' is
\qfullbegin
\qn{15} Number of compounds, which can evolve $\mathrm{CO_2}$ on heating, are $Z$, what is $Z\times5=\ldots\ldots..$
\qfig{p1_chem_s3_q15_1}
\qfullend
\qn{16} An unknown compound having molecular formula $\mathrm{C_5H_{10}O_5}$ was reacted with excess of $\mathrm{CH_3COCl}$ and pyridine to give compound having molecular weight 318 gm. The number of hydroxy groups present in the unknown compound is $Z$, then the value of $(Z\times15)$ is
\qfullbegin
\qn{17} `$x$' number of alcohols among the following will give positive Iodoform test. Then $x\times10=$ ?
\qfig{p1_chem_s3_q17_9}
\qn{18} Number of compounds from below which can reduce Tollen's reagent is $x$. What is $(x)\times11$?
\qfig{p1_chem_s3_q18_1}
\qfullend
\partsec{3}{MATHEMATICS}
\sect{SECTION-I (i)}
\qn{1} If $x$ is real, then the maximum value of $y=2(a-x)\left(x+\sqrt{x^2+b^2}\right)$ is (given $a\ne0$)
\begin{qoptions}
\opt{A}{$a^2+b^2$}
\opt{B}{$a^2$}
\opt{C}{$\sqrt{a^2+b^2}$}
\opt{D}{$a\sqrt{a^2+b^2}$}
\end{qoptions}
\qn{2} If $\tan A$ and $\tan B$ are the roots of the quadratic equation $x^2-px+q=0$, then value of $\sin^2(A+B)$ is
\begin{qoptions}
\opt{A}{$\dfrac{p^2}{p^2+q^2}$}
\opt{B}{$\dfrac{p^2}{(q+p)^2}$}
\opt{C}{$1-\dfrac{p}{(1-q)^2}$}
\opt{D}{$\dfrac{p^2}{(1-q)^2+p^2}$}
\end{qoptions}
\qn{3} For $a,b>0$, the value of {\small\[\int_0^1\dfrac{\log\left(\dfrac{(x+a)^{(x+a)}}{(x+b)^{(x+b)}}\right)}{(x+a)(x+b)\log(x+a)\log(x+b)}\,dx\]} is equal to
\begin{qoptions}
\opt{A}{$\log\!\left(\dfrac{1+a}{1+b}\right)\dfrac{\log a}{\log b}$}
\opt{B}{$\dfrac{\log(1+a)\log a}{\log(1+b)\log b}$}
\opt{C}{$\log\!\left(\dfrac{\log(1+a)\log a}{\log(1+b)\log b}\right)$}
\opt{D}{$\log\!\left(\dfrac{\log(1+b)\log a}{\log(1+a)\log b}\right)$}
\end{qoptions}
\qn{4} Through any point $(x,y)$ of a curve which passes through the origin, lines are drawn parallel to the coordinate axes. The curve, given that it divides the rectangle formed by the two lines and the axes into two areas, one of which is twice the other, represents a family of
\begin{qoptions}
\opt{A}{Circles}
\opt{B}{Parabolas}
\opt{C}{Hyperbolas}
\opt{D}{Straight lines}
\end{qoptions}
\qn{5} Consider the following regions in the plane

$R_1=\{(x,y):0\le x\le1\text{ and }0\le y\le1\}$ \quad $R_2=\{(x,y):x^2+y^2\le4/3\}$

The area of the region $R_1\cap R_2$ can be expressed as $\dfrac{a\sqrt3+b\pi}{9}$, where $a$ and $b$ are integers. Then the value of $(a+b)$ equals
\begin{qoptions}
\opt{A}{2}
\opt{B}{3}
\opt{C}{4}
\opt{D}{5}
\end{qoptions}
\sect{SECTION-I (ii)}
\qn{6} Let $a,b,c,p,q$ be five different non-zero real numbers and $x,y,z$ be three numbers satisfying the system of equations
\[\dfrac xa+\dfrac y{a-p}+\dfrac z{a-q}=1\]
\[\dfrac xb+\dfrac y{b-p}+\dfrac z{b-q}=1\]
\[\dfrac xc+\dfrac y{c-p}+\dfrac z{c-q}=1\]
then
\begin{qoptions}
\opt{A}{$x+y+z=a+b+c-p-q$}
\opt{B}{$x=\dfrac{abc}{pq}$}
\opt{C}{$y=\dfrac{(a-p)(b-p)(c-p)}{p(p-q)}$}
\opt{D}{$z=\dfrac{(a-q)(b-q)(c-q)}{q(q-p)}$}
\end{qoptions}
\qn{7} Suppose $m\in R$. The quadratic equation $x^2-(m-3)x+m=0$ has
\begin{qoptions}
\opt{A}{Real and distinct roots if and only if $m\in(-\infty,1)\cup(9,\infty)$}
\opt{B}{Both positive roots if and only if $m\in[9,\infty)$}
\opt{C}{Both negative roots if and only if $m\in(0,1]$}
\opt{D}{At least one positive root if and only if $m\in(9,\infty)$}
\end{qoptions}
\qn{8} For $k\in N$, let $b_k=\displaystyle\int_1^e\dfrac{1}{(k+\log x)(k+1+\log x)}\dfrac{dx}{x}$ and $a_n=\displaystyle\sum_{k=1}^{n}b_k$. Then
\begin{qoptions}
\opt{A}{$b_k=\log\!\left(\dfrac{(k+1)^2}{k(k+2)}\right)$}
\opt{B}{$a_n=\log\!\left(\dfrac{2n+2}{n+2}\right)$}
\opt{C}{$\displaystyle\lim_{n\to\infty}a_n=\log2$}
\opt{D}{$\displaystyle\lim_{n\to\infty}a_n<1$}
\end{qoptions}
\qn{9} For $a>0$, let $I(a)=\displaystyle\int_0^a\dfrac{dx}{1+e^{a^2-2ax}}$ and let $0<\theta<\dfrac\pi2$, then
\begin{qoptions}
\opt{A}{$I'(a)=\ln\sqrt e$}
\opt{B}{$I(\sin\theta)+I(\cos\theta)\le\sin(\pi/4)$}
\opt{C}{$I(\sin^2\theta)+I(\cos^2\theta)=\dfrac12$}
\opt{D}{$I(2a-a^2)\le\dfrac12$}
\end{qoptions}
\qn{10} If $y=f(x)$ is a solution of $x\dfrac{dy}{dx}-y=x\tan\dfrac yx,\ y(1)=\dfrac\pi2$, then
\begin{qoptions}
\opt{A}{$\sin\dfrac yx=x$}
\opt{B}{$\displaystyle\int_0^1y\,dx=\dfrac38\pi$}
\opt{C}{$\operatorname{cosec}\dfrac yx=|x|$}
\opt{D}{$\displaystyle\int_0^1y\,dx=\dfrac\pi8$}
\end{qoptions}
\qn{11} Let $f$ be a nonnegative function defined on the interval $[0,\pi/2]$. If $\displaystyle\int_0^x\left(f'(t)-\sin2t\right)dt=\int_x^0f(t)\tan t\,dt$ and $f(0)=1$, then
\begin{qoptions}
\opt{A}{$f(0)=1$ is the absolute maximum value of $f$}
\opt{B}{$f(0)=1$ is the minimum value of $f$}
\opt{C}{$\displaystyle\int_0^{\pi/2}y\,dx=3-\dfrac\pi2$}
\opt{D}{$f(\pi/3)=1$}
\end{qoptions}
\qn{12} Area of the region bounded by the curve $y=\tan x$ and lines $y=0$ and $x=1$ is equal to
\begin{qoptions}
\opt{A}{$\displaystyle\int_0^1\tan(1-x)\,dx$}
\opt{B}{$\displaystyle\int_0^{\tan1}\tan^{-1}y\,dy$}
\opt{C}{$\displaystyle\int_0^1\tan^{-1}x\,dx$}
\opt{D}{$\tan1-\displaystyle\int_0^{\tan1}\tan^{-1}x\,dx$}
\end{qoptions}
\qn{13} Let $f:[0,1]\to R$ be continuous and differentiable function such that for any $x,y\in[0,1]$ $xf(y)+yf(x)\le1$, then which of the following options is/are correct?
\begin{qoptions}
\opt{A}{$\displaystyle\int_0^1f(x)\,dx\le\dfrac\pi4$}
\opt{B}{$\displaystyle\int_0^1x^2f'(x)\,dx\ge f(1)-1$}
\opt{C}{$\displaystyle\int_0^1xf'(x)\,dx\ge f(1)-\dfrac\pi4$}
\opt{D}{$\displaystyle\int_0^1f(x)\,dx\le\dfrac\pi8$}
\end{qoptions}
\sect{SECTION-II}
\qn{14} The parabola $P:y=ax^2$ where `$a$' is a positive real constant, is touched by the line $L:y=mx-b$ (where $m$ is a positive constant and $b$ is real) at the point $T$. Let $Q$ be the point of intersection of the line $L$ and the $y$-axis such that $TQ=1$. If $A$ denotes the maximum value of the area of the region surrounded by $P$, $L$ and the $x$-axis, find the value of $\dfrac1A$.
\qn{15} Let $y=f(x)$ be a curve $C_1$ passing through $(2,2)$ and $\left(8,\dfrac12\right)$ and satisfying a differential equation $y\left(\dfrac{d^2y}{dx^2}\right)=2\left(\dfrac{dy}{dx}\right)^2$. Curve $C_2$ is the director circle of the circle $x^2+y^2=2$. If the shortest distance between the curves $C_1$ and $C_2$ is $\left(\sqrt p-q\right)$ where $p,q\in N$, then find the value of $\left(p^2-q\right)$.
\qn{16} Find the number of values of $x$ satisfying $\displaystyle\int_0^xt^2\sin(x-t)\,dt=x^2$ in $[0,100]$.
\qn{17} Let $P(x)$ be a polynomial with real coefficients such that $\left(x^2+x+1\right)P(x-1)=\left(x^2-x+1\right)P(x)\ \forall x\in R$ and $P(1)=3$.

If $\displaystyle\int_0^1\tan^{-1}\!\left(\dfrac{2x}{1+P\left(x^2\right)}\right)dx+\int_0^1\tan^{-1}(x+1)\,dx=\dfrac{k}{16}(\pi-\ln4)$, then find the value of $k$.
\qn{18} Sum of all integral value(s) of $k$ for which $|x-1|+\left|x^2-4\right|+|x+1|=k$ has 2 distinct real solutions in the interval $x\in[-3,4]$ is
\papersec{2}
\partsec{1}{PHYSICS}
\sect{SECTION-I}
\qn{1} The current in a long solenoid is increasing linearly with time, so the flux is proportional to $t:\Phi=\alpha t$. Two voltmeters are connected to diametrically opposite points ($A$ and $B$), together with resistors $\left(R_1\text{ and }R_2\right)$, as shown in fig. Assume that these are ideal voltmeters that draw negligible current (they have huge internal resistance), and that a voltmeter registers $-\displaystyle\int_a^bE\cdot d\ell$ between the terminals and through the meter.
\qfig{p2_phy_s1_q01_1}
\begin{qoptions}
\opt{A}{$V_1=\dfrac{\alpha R_1}{\left(R_1+R_2\right)}$}
\opt{B}{$V_2=-\dfrac{\alpha R_2}{\left(R_1+R_2\right)}$}
\opt{C}{$V_1=V_2=\alpha$}
\opt{D}{$V_1=V_2=\dfrac{\alpha R_1R_2}{R_1+R_2}$}
\end{qoptions}
\qn{2} A circular wire loop (radius $r$, resistance $R$) encloses a region of uniform magnetic field, $B$, perpendicular to its plane. The field increases linearly with time $(B=\alpha t)$. An ideal voltmeter (infinite internal resistance) is connected between points $P$ and $Q$.
\qfig{p2_phy_s1_q02_1}
\begin{qoptions}
\opt{A}{The current in the loop is $\dfrac{\pi\alpha r^2}{2R}$}
\opt{B}{The reading of voltmeter is $\dfrac{\alpha r^2}{2}$}
\opt{C}{The reading of voltmeter is $\alpha r^2$}
\opt{D}{The current in loop is $\dfrac{\pi\alpha r^2}{R}$}
\end{qoptions}
\qn{3} A conducting rod of length $\ell$ slides in a vertical plane against two insulated perpendicular surfaces as shown. Find the induced emf across the rod if its end $A$ is pulled with $\vec v=v_0\hat i$. Outward uniform magnetic field of induction $B$ is applied.
\qfig{p2_phy_s1_q03_1}
\begin{qoptions}
\opt{A}{$\varepsilon=\dfrac{B\ell V_0}{2}$ when $\theta=30^\circ$}
\opt{B}{$\varepsilon=B\ell V_0$ when $\theta=30^\circ$}
\opt{C}{$\varepsilon=0$ when $\theta=45^\circ$}
\opt{D}{$\varepsilon=\dfrac{B\ell V_0}{\sqrt2}$ when $\theta=45^\circ$}
\end{qoptions}
\qn{4} Given figure is a device to fire (project) a body (sliding bar) with a great speed in the order of km/s by an Ampere's force on the conducting sliding bar offered by an impulsive current sent by the charged capacitor after closing the key in the inward magnetic field.

If $Q_0$ is the initial charge of the capacitor $C$, $M=$ mass of the sliding bar and $B=$ magnetic field (produced by the closely spaced parallel rails assumed to be uniform), answer the following questions: (at $t=0$, $K$ is closed)
\qfig{p2_phy_s1_q04_1}
\begin{qoptions}
\opt{A}{The final charge on capacitor $\dfrac{B^2\ell^2CQ_0}{M+CB^2\ell^2}$}
\opt{B}{The maximum (muzzle) speed of the projectile is $\dfrac{2B\ell Q_0}{M+CB^2\ell^2}$}
\opt{C}{The maximum (muzzle) speed of the projectile is $\dfrac{B\ell Q_0}{M+CB^2\ell^2}$}
\opt{D}{The final charge on capacitor $\dfrac{2B^2\ell^2CQ_0}{M+CB^2\ell^2}$}
\end{qoptions}
\qn{5} A small loop of radius $r(\ll R)$, where $R=$ radius of the larger loop that carries a current $i$ is placed on the axis of the bigger loop. If the smaller loop is rotated through $90^\circ$, then the (the resistance of the small loop is $R_0$). [$\omega=$ Angular velocity of loop along the diametrical axis]
\qfig{p2_phy_s1_q05_1}
\begin{qoptions}
\opt{A}{Change in flux is $\dfrac{\pi\mu_0r^2i}{4\sqrt2R}$}
\opt{B}{Average current flowing through smaller loop is $\dfrac{\mu_0r^2\omega i}{2\sqrt2R_0R}$}
\opt{C}{Change in flux is $\dfrac{\pi\mu_0r^2i}{2\sqrt2R}$}
\opt{D}{Average current flowing through smaller loop is $\dfrac{\mu_0r^2\omega i}{\sqrt2R_0R}$}
\end{qoptions}
\qn{6} \qfigi{p2_phy_s1_q05_2}

$R=100\,\Omega,\ L=\dfrac2\pi$ H, $C=\dfrac{1}{10\pi}$ mF

In above given AC circuit, $V_R,V_L$ and $V_C$ are potential drop across resistance, inductor and capacitor at an instant $t$. Then
\begin{qoptions}
\opt{A}{$V_C=220\sin(100\pi t-\pi/4)$}
\opt{B}{$V_L=440\sin(100\pi t+\pi/4)$}
\opt{C}{$V_C=220\sin(100\pi t-3\pi/4)$}
\opt{D}{$V_L=220\sin(100\pi t-\pi/4)$}
\end{qoptions}
\qn{7} The potential energy between two atoms in a diatomic molecule varies with $x$ as $U=\dfrac{a}{x^{12}}-\dfrac{b}{x^6}$, where $a$ and $b$ are positive constants. Consider oscillation of one atom if the other atom is kept stationary.
\qfig{p2_phy_s1_q07_1}
\begin{qoptions}
\opt{A}{Equivalent spring constant $=\left(\dfrac{9b^2}{a}\right)\left(\dfrac{b}{2a}\right)^{1/3}$}
\opt{B}{Equivalent spring constant $=\left(\dfrac{9b^2}{a}\right)\left(\dfrac{2b}{a}\right)^{1/3}$}
\opt{C}{Equilibrium position of particle $=\left(\dfrac{2a}{b}\right)^{1/6}$}
\opt{D}{Equivalent position of particle $=\left(\dfrac{a}{2b}\right)^{1/6}$}
\end{qoptions}
\qn{8} A uniform disc of mass $m$ and radius $R$ is hanging from a rigid support by $N$-number of strings each of length $h$ as shown in the figure. Find the time period of small torsional oscillation about vertical axis passing through the CM of the disc.
\qfig{p2_phy_s1_q08_1}
\begin{qoptions}
\opt{A}{$\pi\left(\dfrac Rr\right)\sqrt{\dfrac hg}$}
\opt{B}{$2\pi\left(\dfrac Rr\right)\sqrt{\dfrac{2h}{g}}$}
\opt{C}{$2\pi(R/r)\sqrt{\dfrac hg}$}
\opt{D}{$\pi(R/r)\sqrt{\dfrac{2h}{g}}$}
\end{qoptions}
\fwstart
\sect{SECTION-IV}
\qn{9} In column -- I time varying magnetic field is available in circular region of radius $R$. $R=1$ m, $\dfrac{dB}{dt}=2$ T/s

In column -- II induced emf or $\oint\vec E\cdot d\vec\ell$ value is given. Match $\varepsilon$ value in column -- II with different situations in column -- I $(\omega=2$ rad/sec and $B=1$ T$)$.\par\nobreak
\begin{center}\begin{tabular}{|L{0.6cm}|L{6.8cm}|L{0.8cm}|L{4.2cm}|}\hline
\multicolumn{2}{|c|}{\textbf{Column I} (Conductor in a time varying magnetic field)} & \multicolumn{2}{c|}{\textbf{Column II} (Induced emf)}\\\hline
(A) & \qfigi[0.8]{p2_phy_s2_q09_A}

$\varepsilon=$ P.D. between points 1 and 2 & (P) & $\varepsilon=3.14$ volt\\\hline
(B) & \qfigi[0.8]{p2_phy_s2_q09_B}

$\varepsilon=$ P.D. between points 1 and 2 & (Q) & $\varepsilon=0.57$ volt\\\hline
(C) & \qfigi[0.8]{p2_phy_s2_q09_3} & (R) & $\varepsilon=3/4$ volt\\\hline
(D) & \qfigi[0.8]{p2_phy_s2_q09_D}

$\varepsilon=\displaystyle\int\vec E\cdot d\vec\ell$ in circular path 132 of radius $2R$ & (S) & $\varepsilon=1$ volt\\\hline
\end{tabular}\end{center}
\fwend
\fwstart
\qn{10} In column -- I AC circuits are given and in column -- II variation of $i$ with respect to $\omega$ is give. Match the following. Value of $R,L$ and $C$ is constant and only $\omega$ of source is varied.\par\nobreak
\begin{center}\begin{tabular}{|L{0.6cm}|L{5.2cm}|L{0.8cm}|L{5.2cm}|}\hline
\multicolumn{2}{|c|}{\textbf{Column I} (A.C. Circuit)} & \multicolumn{2}{c|}{\textbf{Column II} $i=f(\omega)$}\\\hline
(A) & \qfigi[0.8]{p2_phy_s2_q10_A} & (P) & \qfigi[0.8]{p2_phy_s2_q10_2}\\\hline
(B) & \qfigi[0.8]{p2_phy_s2_q10_B} & (Q) & \qfigi[0.8]{p2_phy_s2_q10_3}\\\hline
(C) & \qfigi[0.8]{p2_phy_s2_q10_C} & (R) & \qfigi[0.8]{p2_phy_s2_q10_5}\\\hline
(D) & \qfigi[0.8]{p2_phy_s2_q10_D} & (S) & \qfigi[0.8]{p2_phy_s2_q10_7}\\\hline
 & & (T) & \qfigi[0.8]{p2_phy_s2_q10_9}\\\hline
\end{tabular}\end{center}
\fwend
\sect{SECTION-II}
\qn{11} A square loop of side $\ell=10$ cm with its sides parallel to $x$ and $y$ axes moves with velocity $\vec v=-7\hat i$ m/s. An outward magentic field is present which decreases along negative $x$--axis with a gradient $\dfrac{\partial B}{\partial x}=0.1$ T/m and increases with time at a rate of $\dfrac{\partial B}{\partial t}=0.1$ T/s. If the resistance of the loop is $2\times10^{-3}$ ohms, then find the magnitude and direction of induced current in the loop. (in Ampere)
\qfig{p2_phy_s3_q11_1}
\qn{12} Two thin long conducting tapes of which width ``$a$'' are separated through a distance ``$h$'' as shown in the figure. Current $a$ is uniformly distributed over (the area of) each tape. Assuming linear density of the current over its width uniform the inductance of the tape per its unit length is $\dfrac{\mu_0h^x}{a^y}$. Find value of $x+y$.
\qfig{p2_phy_s3_q12_1}
\qn{13} A resistor $R$ is connected in series with a coil (Resistance not negligible). The system is subjected to an A.C. Supply of peak voltage $V_0$. If the peak voltages across the resistor $R$ and the coils are $V_1$ and $V_2$ respectively, find the power (in watt) dissipated in the coil. (assume $V_1=5$V, $V_2=7$V, $V_0=10$V, $R=13\,\Omega$)
\qn{14} In $R$--$L$--$C$ series circuit, we have same current at angular frequencies $\omega_1$ and $\omega_2$. Find the resonant frequency. (Take $\omega_1=4$ units, $\omega_2=9$ units)
\qn{15} A body of mass $m$ hangs from a smooth fixed pully $P_1$ by the inextensible string fitted with the springs of stiffnesses $k_1$ and $k_3$. The string passes over the smooth light pulley $P_2$ which is connected with another ideal spring of stiffness $k_2$. Find the period of oscillation of the body. (in seconds) assume all springs masslass and take $k_1=k_3=\dfrac{k_2}{4}=\pi^2$ SI units, $M=3$ kg
\qfig{p2_phy_s3_q15_1}
\qn{16} A stepped pully having mass $m$, radius $R$ and radius of gyration about cm is connected with two ideal springs of stiffness $k_1$ and $k_2$ as shown in the figure. If the pulley rolls without sliding, the angular frequency of its oscillation is given by $\sqrt{\dfrac{xk_1r^2+yk_2(R-r)^2}{M\left(K^2+r^2\right)}}$. Find value of $x+y$.
\qfig{p2_phy_s3_q16_2}
\qn{17} A simple pendulum has time period $T_0$, say. When it is placed in a liquid the time period decreases. Then we need to move the point of suspension with a vertical acceleration $a$ $(a=2g)$ to regain the time period of simple pendulum. The ratio of densities of liquid and the bob is:
\qfig{p2_phy_s3_q17_1}
\qn{18} Period of oscillation of the particle $P$ is $2\pi\sqrt{\dfrac{3m}{xk}}$, then $x$ is equal to
\qfig{p2_phy_s3_q18_1}
\partsec{2}{CHEMISTRY}
\sect{SECTION-I}
\qn{1} A reversible cyclic process for an ideal gas is shown below. Here, $p$, $V$ and $T$ are pressure, volume and temperature, respectively. The thermodynamic parameters $q$, $w$, $H$ and $U$ are heat, work, enthalpy and internal energy, respectively.
\qfig{p2_chem_s1_q01_1}
The correct options is (are)
\begin{qoptions}
\opt{A}{$q_{AC}=\Delta U_{BC}$ and $W_{AB}=nRT_1\ell n\dfrac{V_1}{V_2}$}
\opt{B}{$W_{BC}=P_2\left(V_1-V_2\right)$ and $W_{AC}=0$}
\opt{C}{$\Delta H_{CA}<\Delta U_{AC}$ and $q_{BC}=\Delta H_{AC}$}
\opt{D}{$q_{AC}=\Delta U_{BC}$ and $\Delta H_{CA}>\Delta U_{CA}$}
\end{qoptions}
\qn{2} No.\,of compound which are better hydride donor than \qfigi{p2_chem_s1_q02_1}
\begin{qoptions}
\opt{A}{\qfigi{p2_chem_s1_q02_A}}
\opt{B}{\qfigi{p2_chem_s1_q02_B}}
\opt{C}{\qfigi{p2_chem_s1_q02_C}}
\opt{D}{\qfigi{p2_chem_s1_q02_D}}
\end{qoptions}
\qn{3} Find out the possible products in the following reaction.
\qfig{p2_chem_s1_q03_1}
\begin{qoptions}
\opt{A}{\qfigi{p2_chem_s1_q03_A}}
\opt{B}{\qfigi{p2_chem_s1_q03_B}}
\opt{C}{$\mathrm{CH_3CH_2OH}$}
\opt{D}{\qfigi{p2_chem_s1_q03_D}}
\end{qoptions}
\qfullbegin
\qn{4} Which of the following reactions represent the major product?
\begin{qoptions}
\opt{A}{\qfigi{p2_chem_s1_q04_A}}
\opt{B}{$\mathrm{Ph-\overset{O}{\overset{\|}{C}}-NH_2}\ \xrightarrow[\mathrm{KOH}]{\mathrm{Br_2}}\ \mathrm{Ph-NH_2}$}
\opt{C}{\qfigi{p2_chem_s1_q04_C}}
\opt{D}{\qfigi{p2_chem_s1_q04_D}}
\end{qoptions}
\qfullend
\qn{5} Which of the following reactants on reaction with conc. NaOH followed by acidification gives following lactone as the main product?
\qfig{p2_chem_s1_q05_1}
\begin{qoptions}
\opt{A}{\qfigi{p2_chem_s1_q05_A}}
\opt{B}{\qfigi{p2_chem_s1_q05_B}}
\opt{C}{\qfigi{p2_chem_s1_q05_C}}
\opt{D}{\qfigi{p2_chem_s1_q05_D}}
\end{qoptions}
\qn{6} Which of the following products is/are correctly mentioned in the following reactions.
\begin{qoptions}
\opt{A}{$\mathrm{HCOH}\xrightarrow{\mathrm{NaOD}}\mathrm{HCOONa}+\mathrm{CH_3OD}$}
\opt{B}{$\mathrm{HCDO}\xrightarrow{\mathrm{NaOH}}\mathrm{DCOONa}+\mathrm{CH_2DOH}$}
\opt{C}{$\mathrm{HCDO}\xrightarrow{\mathrm{NaOD}}\mathrm{DCOONa}+\mathrm{DCH_2OD}$}
\opt{D}{$\mathrm{D_2CO}\xrightarrow{\mathrm{NaOH}}\mathrm{DCOONa}+\mathrm{CD_3OD}$}
\end{qoptions}
\qn{7} $\mathrm{C_4H_6O_4\,(X)\xrightarrow{LAH}A\xrightarrow{PCl_3}B\xrightarrow{NaCN}C}$ $\mathrm{\xrightarrow{H_3O^+,\,CaO,\,\Delta}Y\xleftarrow{CaO,\,\Delta}}$ hexanedioic acid
\begin{qoptions}
\opt{A}{$B$ is \begin{tabular}[t]{@{}l@{}}$\mathrm{CH_2-CH_2-Cl}$\\$|$\\$\mathrm{CH_2-CH_2-Cl}$\end{tabular}}
\opt{B}{$Y$ is \qfigi{p2_chem_s1_q07_1}}
\opt{C}{$X$ is \begin{tabular}[t]{@{}l@{}}$\mathrm{CH_2-COOH}$\\$|$\\$\mathrm{CH_2-COOH}$\end{tabular}}
\opt{D}{$C$ is \begin{tabular}[t]{@{}l@{}}$\mathrm{CH_2-CH_2-COOH}$\\$|$\\$\mathrm{CH_2-CH_2-COOH}$\end{tabular}}
\end{qoptions}
\qn{8} A sample containing 1.0 mol of an ideal gas is expanded isothermally and reversibly to ten times of its original volume, in two separate experiments. The expansion is carried out at 300 K and at 600 K respectively. Choose the correct option.
\begin{qoptions}
\opt{A}{Work done at 600 K is 20 times the work done at 300 K}
\opt{B}{Work done at 300 K is twice the work done at 600 K}
\opt{C}{Work done at 600 K is twice the work done at 300 K}
\opt{D}{$\Delta U=0$ in both cases}
\end{qoptions}
\fwstart
\sect{SECTION-IV}
\qn{9} Match the compounds/ions in Column I with their properties/reactions in Column II.\par\nobreak
\begin{center}\begin{tabular}{|L{0.6cm}|L{5.2cm}|L{0.8cm}|L{6.4cm}|}\hline
\multicolumn{2}{|c|}{\textbf{Column I}} & \multicolumn{2}{c|}{\textbf{Column II}}\\\hline
(A) & \qfigi[0.8]{p2_chem_s2_q09_A} & (P) & Gives red precipitate with $\mathrm{CuSO_4}$ / sodium potassium tartrate solution\\\hline
(B) & \qfigi[0.8]{p2_chem_s2_q09_B} & (Q) & Shows tautomerism\\\hline
(C) & \qfigi[0.8]{p2_chem_s2_q09_C} & (R) & Gives two isomeric compounds with hydrazine under acidic condition\\\hline
(D) & \qfigi[0.8]{p2_chem_s2_q09_D} & (S) & Gives yellow precipitate with $\mathrm{I_2}$ / NaOH\\\hline
 & & (T) & Gives silver mirror test with tollens reagent\\\hline
\end{tabular}\end{center}
\fwend
\fwstart
\qn{10} Match the following\par\nobreak
\begin{center}\begin{tabular}{|L{0.6cm}|L{5.6cm}|L{0.8cm}|L{5.2cm}|}\hline
\multicolumn{2}{|c|}{\textbf{Column I}} & \multicolumn{2}{c|}{\textbf{Column II}}\\\hline
(A) & Isothermal reversible expansion & (P) & $\Delta S_{\mathrm{system}}>0$\\\hline
(B) & Isothermal irreversible & (Q) & $\Delta S_{\mathrm{system}}=0$\\\hline
(C) & Adiabatic reversible expansion & (R) & $\Delta S_{\mathrm{universe}}>0$\\\hline
(D) & Adiabatic irreversible expansion & (S) & $\Delta S_{\mathrm{universe}}=0$\\\hline
 & & (T) & $\Delta S_{\mathrm{surrounding}}<0$\\\hline
\end{tabular}\end{center}
\fwend
\sect{SECTION-II}
\qn{11} The number of products possible for $x$ (including stereomers) in the following reaction is:
\qfig{p2_chem_s3_q11_1}
\qn{12} Which of the reactants have maximum rate of reaction as compared to Benzaldehyde for the formation of cyanohydrin compound

1. p--Methyoxy benzaldehyde \quad 2. p--Nitrobenzaldehyde \quad 3. p--Methylbenzaldehyde \quad 4. p--aminobenzaldehyde \quad 5. p--Hydroxybenzaldehyde \quad 6. p--Anilinium benzaldehyde
\qfullbegin
\qn{13} How many of the following compounds can undergo decarboxylation on heating.
\qfig{p2_chem_s3_q13_10}
\qfullend
\qn{14} An organic compound `A' on treatment with ethyl alcohol gives a carboxylic acid `B' and compound `C'. Hydrolysis of `C' under acidic conditions gives `B' and `D'. Oxidation of `D' with $\mathrm{KMnO_4}$ also gives `B'. `B' on heating with $\mathrm{Ca(OH)_2}$ gives `E' $\left(\mathrm{C_3H_6O}\right)$. E does not give Tollen's test and does not reduce Fehling's solution but form a 2,4--dinitrophenyl hydrazone. Let no.\,of carbons are present in compound `A' is `x' and oxygen be `y', then $x+y=$ ?
\qn{15} Number of reactions in which product has one carbon lesser than reactant

\qfig{p2_chem_s3_q15_1}
\noindent(II) $\mathrm{CH_3CONH_2+OBr^-\xrightarrow[\Delta]{H_3O^\oplus}}$

\noindent(III) $\mathrm{CH_3COOH+N_3H+H_2SO_4\xrightarrow{\Delta}}$
\qfig{p2_chem_s3_q15_2}
\qfig{p2_chem_s3_q15_3}
\qfullbegin
\qn{16} Which of the following will show `+' ve haloform test?
\qfig{p2_chem_s3_q16_9}
\qfullend
\qn{17} Isothermal compressibility $(K)$ of an ideal gas is defined as $K=-\dfrac1V\left(\dfrac{\partial V}{\partial P}\right)_{T,n}$. What is the isothermal compressibility factor (per atm) for an ideal gas at 1.0 atm?
\qn{18} Find the total number of products obtained when
\qfig{p2_chem_s3_q18_1}
Undergoes decarboxylation
\partsec{3}{MATHEMATICS}
\sect{SECTION-I}
\qn{1} If $f(x)=\displaystyle\lim_{m\to\infty}\dfrac1m\sum_{k=0}^{m}\left|x-\dfrac km\right|$; then
\begin{qoptions}
\opt{A}{$f(0)=\dfrac12$}
\opt{B}{$f(1)=\dfrac12$}
\opt{C}{$f(-1)=\dfrac32$}
\opt{D}{$f(2)=\dfrac32$}
\end{qoptions}
\qn{2} Let $f(x)$ is a differentiable function such that $f'(x)=f(x)+\displaystyle\int_0^2f(x)\,dx$ if $f(0)=\dfrac{4-e^2}{3}$, then
\begin{qoptions}
\opt{A}{$f(x)=e^x-\left(\dfrac{e^2-1}{3}\right)$}
\opt{B}{$f(x)=e^x+\left(\dfrac{e^2-1}{3}\right)$}
\opt{C}{$f(0)<0$}
\opt{D}{$f(x)$ has atleast one root between 0 \& 1}
\end{qoptions}
\qn{3} The differential equation $\left\{xy^3(1+x\cos x)-y\right\}dx+x\,dy=0$ represents the curve $\dfrac{x^2}{ay^2}=\dfrac{x^3}{b}+x^2\sin x+cx\cos x-d\sin x+k$. Then
\begin{qoptions}
\opt{A}{$c+d=0$}
\opt{B}{$a+b+c+d=9$}
\opt{C}{$b+d=-1$}
\opt{D}{$a+b=5$}
\end{qoptions}
\qn{4} Solution of equation $y\,dx\left(1+x^2\cdot\ln xy\right)+x\,dy\left(1+y^2\ln(xy)\right)=0$ is
\begin{qoptions}
\opt{A}{$\ln(\ln xy)+x^2+y^2+c=0$}
\opt{B}{$\ln xy=ce^{\frac{\left(x^2+y^2\right)}{2}}$}
\opt{C}{$2\ln(\ln xy)+x^2+y^2+c=0$}
\opt{D}{$\ln xy=ce^{-\frac{\left(x^2+y^2\right)}{2}}$}
\end{qoptions}
\qn{5} If $ax^2+bx+c=0$ and $cx^2+bx+a=0$ $(a,b,c\in R)$ have a common non-real root, then
\begin{qoptions}
\opt{A}{$-2|a|<b<2|a|$}
\opt{B}{$-2|c|<|b|<2|c|$}
\opt{C}{$a=\pm c$}
\opt{D}{$a=c$}
\end{qoptions}
\qn{6} The equation $a_8x^8+a_7x^7+a_6x^6+\ldots..+a_0=0$ has all its roots positive and real (where $a_8=1,a_7=-4,a_0=1/2^8$), then
\begin{qoptions}
\opt{A}{$a_1=\dfrac1{2^8}$}
\opt{B}{$a_1=-\dfrac1{2^4}$}
\opt{C}{$a_2=\dfrac7{2^4}$}
\opt{D}{$a_3=\dfrac1{2^2}$}
\end{qoptions}
\qn{7} A differentiable function satisfies $f(x)=\displaystyle\int_0^x\left\{f(t)\cos t-\cos(t-x)\right\}dt$. Which of the following hold good?
\begin{qoptions}
\opt{A}{$f(x)$ has a minimum value $1-e$}
\opt{B}{$f(x)$ has a maximum value $1-e^{-1}$}
\opt{C}{$f''\!\left(\dfrac\pi2\right)=e$}
\opt{D}{$f'(0)=1$}
\end{qoptions}
\qn{8} A curve $y=f(x)$ passes through $(1,1)$ and tangent at $P(x,y)$ cuts the $x$-axis at $A$ and $B$ respectively such that $BP:AP=3:1$, then
\begin{qoptions}
\opt{A}{Equation of curve is $x^3y+1=0$}
\opt{B}{Tangent at $(1,1)$ is $x+3y=4$}
\opt{C}{Curve passes through $\left(3,\dfrac1{27}\right)$}
\opt{D}{Equation of curve is $x^3y=1$}
\end{qoptions}
\fwstart
\sect{SECTION-IV}
\qn{9} Let $a<b<c<d$, match with the column in which equation has root in:\par\nobreak
\begin{center}\begin{tabular}{|L{0.6cm}|L{6.6cm}|L{0.8cm}|L{3.6cm}|}\hline
\multicolumn{2}{|c|}{\textbf{Column I}} & \multicolumn{2}{c|}{\textbf{Column II}}\\\hline
(A) & $(x-a)(x-c)+2007(x-b)(x-d)=0$ & (P) & $(c,d)$\\\hline
(B) & $2007(x-a)(x-c)-2006(x-b)(x-d)=0$ & (Q) & $(a,b)$\\\hline
(C) & $(x-a)(x-b)(x-c)+(x-d)=0$ & (R) & $(-\infty,a)$\\\hline
(D) & $(x-a)(x-b)+(x-c)=0$ & (S) & $(b,c)$\\\hline
\end{tabular}\end{center}
\fwend
\fwstart
\qn{10} Match curves with their Differential Equations:\par\nobreak
\begin{center}\begin{tabular}{|L{0.6cm}|L{5.0cm}|L{0.8cm}|L{6.4cm}|}\hline
\multicolumn{2}{|c|}{\textbf{Column I}} & \multicolumn{2}{c|}{\textbf{Column II}}\\\hline
(A) & $y=c_1e^{2x}+c_2e^{3x}+c_3e^{-x}$ & (P) & $\dfrac{d^3y}{dx^3}-6\dfrac{d^2y}{dx^2}+12\dfrac{dy}{dx}-8y=0$\\\hline
(B) & $y=\left(c_1+c_2x+c_3x^2\right)e^{2x}$ & (Q) & $x\left[y\dfrac{d^2y}{dx^2}+\left(\dfrac{dy}{dx}\right)^2\right]=y\dfrac{dy}{dx}$\\\hline
(C) & $xy=ae^x+be^{-x}+x^2$ & (R) & $\dfrac{d^3y}{dx^3}-4\dfrac{d^2y}{dx^2}+\dfrac{dy}{dx}+6y=0$\\\hline
(D) & $ax^2+by^2=1$ & (S) & $x\dfrac{d^2y}{dx^2}+2\dfrac{dy}{dx}-xy+x^2-2=0$\\\hline
\end{tabular}\end{center}
\fwend
\sect{SECTION-II}
\qn{11} The area of the region bounded by the curves $2^{|x|}|y|+2^{|x|-1}\le1,\ |x|\le\dfrac12$ and $|y|\le\dfrac12$ is $\dfrac{a}{\log2}\left(1-\dfrac1{\sqrt b}\right)+c$ then $a+b+c$ is $(a,b,c\in N)$
\qn{12} The area bounded by $y=x^2+1$ and the tangents to it drawn from the origin is $A$ then $3A$ is
\qn{13} If {\small\[\begin{aligned}&\int_{-5}^{-2}\left(\dfrac{x^2-x}{x^3-3x+1}\right)^2dx+\int_{1/6}^{1/3}\left(\dfrac{x^2-x}{x^3-3x+1}\right)^2dx\\&+\int_{6/5}^{3/2}\left(\dfrac{x^2-x}{x^3-3x+1}\right)^2dx=\dfrac pq\end{aligned}\]} where $p$ and $q$ are co--primes, then $\dfrac{(6q-p)}{6}$ is
\qn{14} The area given by $x+y\le6,\ x^2+y^2\le6y$ and $y^2\le8x$ is $\dfrac ab\pi-\dfrac16$ then $a+b$ is (given $a,b\in N$ $a$, $b$ are relatively prime)
\qn{15} Let $P(x)=x^6-x^5-x^3-x^2-x$ and $\alpha,\beta,\gamma,\delta$ are the roots of the equation $x^4-x^3-x^2-1=0$, then $P(\alpha)+P(\beta)+P(\gamma)+P(\delta)=$
\qn{16} $\displaystyle\int_0^{2\pi}\max\left\{\sin x,\sin^{-1}(\sin x)\right\}dx=\dfrac{\pi^b}{a}-c$, then $(a+b+c)$ is $(a,b,c\in N)$
\qn{17} If real roots of $(x+3)^4+(x+5)^4=16$ are $a$ \& $b$ then $\dfrac{a^2+b^2+ab}{7}$ is
\qn{18} The number of integral values of $a$ for which the inequality $3-|x-a|>x^2$ is satisfied by at least one negative $x$, must be equal to
\end{multicols}
\end{document}
